\documentclass[aps,preprint]{revtex4}%
\usepackage{amsfonts}
\usepackage{amsmath}
\usepackage{amssymb}
\usepackage{graphicx}%
\setcounter{MaxMatrixCols}{30}
%TCIDATA{OutputFilter=latex2.dll}
%TCIDATA{Version=4.10.0.2363}
%TCIDATA{CSTFile=revtex4.cst}
%TCIDATA{Created=Saturday, July 31, 2004 09:56:36}
%TCIDATA{LastRevised=Thursday, August 19, 2004 13:33:57}
%TCIDATA{<META NAME="GraphicsSave" CONTENT="32">}
%TCIDATA{<META NAME="DocumentShell" CONTENT="Articles\SW\REVTeX 4">}
\newtheorem{theorem}{Theorem}
\newtheorem{acknowledgement}[theorem]{Acknowledgement}
\newtheorem{algorithm}[theorem]{Algorithm}
\newtheorem{axiom}[theorem]{Axiom}
\newtheorem{claim}[theorem]{Claim}
\newtheorem{conclusion}[theorem]{Conclusion}
\newtheorem{condition}[theorem]{Condition}
\newtheorem{conjecture}[theorem]{Conjecture}
\newtheorem{corollary}[theorem]{Corollary}
\newtheorem{criterion}[theorem]{Criterion}
\newtheorem{definition}[theorem]{Definition}
\newtheorem{example}[theorem]{Example}
\newtheorem{exercise}[theorem]{Exercise}
\newtheorem{lemma}[theorem]{Lemma}
\newtheorem{notation}[theorem]{Notation}
\newtheorem{problem}[theorem]{Problem}
\newtheorem{proposition}[theorem]{Proposition}
\newtheorem{remark}[theorem]{Remark}
\newtheorem{solution}[theorem]{Solution}
\newtheorem{summary}[theorem]{Summary}
\newenvironment{proof}[1][Proof]{\noindent\textbf{#1.} }{\ \rule{0.5em}{0.5em}}
\begin{document}
\preprint{HEP/123-qed}
\title[Short title for running header]{Long title that appears on the title page}
\author{First Author}
\affiliation{My Institution}
\author{Second Author}
\affiliation{My Institution}
\author{Third Author}
\affiliation{Other Institution}
\keywords{one two three}
\pacs{PACS number}

\begin{abstract}
Shell document for REV\TeX{} 4.

\end{abstract}
\volumeyear{year}
\volumenumber{number}
\issuenumber{number}
\eid{identifier}
\date[Date text]{date}
\received[Received text]{date}

\revised[Revised text]{date}

\accepted[Accepted text]{date}

\published[Published text]{date}

\startpage{101}
\endpage{102}
\maketitle
\tableofcontents


The relaxation of the occupation numbers from integer values and the
determination of optimal phases reflect the correlation effects.%

\begin{equation}
\mathcal{E}\left(  D^{2N}\right)  =\frac{1}{\operatorname{Tr}\left\{
D^{2N}\right\}  }\operatorname{Tr}\left\{  \left\{  \sum_{1\leq i,k\leq
2s}h_{1ik}a_{i}^{\dagger}a_{k}+\frac{1}{4}\sum_{1\leq i,j,k,l\leq
2s}\left\langle ij||kl\right\rangle a_{i}^{\dagger}a_{j}^{\dagger}a_{l}%
a_{k}\right\}  D^{2N}\right\}  , \label{en1}%
\end{equation}%
\begin{equation}
\mathcal{E}\left(  \mathbb{D}^{2}\right)  =\frac{1}{\mathcal{S}}\left\{
\operatorname{Tr}\left\{  \mathbb{H}_{1}\mathbb{D}^{1}\right\}  +\frac{1}%
{4}\operatorname{Tr}\left\{  \mathbb{V}_{2}\mathbb{D}^{2}\right\}  \right\}
\end{equation}%
\begin{equation}
\mathcal{E}\left(  \mathbb{D}^{2}\right)  =\frac{1}{\mathcal{S}}\left\{
\operatorname{Tr}\left\{  \mathbb{H}_{1}\mathbb{D}^{1}\right\}  +\frac{1}%
{4}\operatorname{Tr}\left\{  \mathbb{V}_{2}\mathbb{D}^{2}\right\}  \right\}
\end{equation}%
\begin{equation}
\mathcal{E}\left(  \mathbb{D}^{2}\right)  =\frac{1}{\mathcal{S}\left(
\mathbb{W\left(  \mathbf{\alpha}\right)  }\right)  }\left\{  \operatorname{Tr}%
\left\{  \mathbb{H}_{1}\mathbb{W\left(  \mathbf{\alpha}\right)  D}_{0}%
^{1}\mathbb{W\left(  \mathbf{\alpha}\right)  }^{\dagger}\right\}  +\frac{1}%
{4}\operatorname{Tr}\left\{  \mathbb{V}_{2}\mathbb{W\left(  \mathbf{\alpha
}\right)  \otimes W\left(  \mathbf{\alpha}\right)  D}_{0}^{2}\mathbb{W\left(
\mathbf{\alpha}\right)  }^{\dagger}\otimes\mathbb{W\left(  \mathbf{\alpha
}\right)  }^{\dagger}\right\}  \right\}
\end{equation}%
\begin{align}
D^{1}  &  =\sum D_{ij}^{1}\left\vert \psi_{i}\right\rangle \left\langle
\psi_{j}\right\vert \\
D^{1}\left(  \mathbf{\alpha}\right)   &  =\sum D_{_{0}ij}^{1}V\left(
\mathbf{\alpha}\right)  \left\vert \psi_{i}\right\rangle \left\langle \psi
_{j}\right\vert V\left(  \mathbf{\alpha}\right)  ^{\dagger}=\sum D_{_{0}%
ij}^{1}\left\vert \psi_{i}\left(  \mathbf{\alpha}\right)  \right\rangle
\left\langle \psi_{j}\left(  \mathbf{\alpha}\right)  \right\vert \\
&  \sum V\left(  \mathbf{\alpha}\right)  _{i^{\prime}i}^{t}D_{_{0}ij}%
^{1}V\left(  \mathbf{\alpha}\right)  _{jj^{\prime}}^{\ast}\left\vert \psi
_{i}\right\rangle \left\langle \psi_{j}\right\vert =\sum W\left(
\mathbf{\alpha}\right)  _{i^{\prime}i}D_{_{0}ij}^{1}W\left(  \mathbf{\alpha
}\right)  _{jj^{\prime}}^{\dagger}\left\vert \psi_{i}\right\rangle
\left\langle \psi_{j}\right\vert
\end{align}%
\begin{equation}
D^{2}\left(  \mathbf{\alpha}\right)  =\frac{1}{4}\sum D_{0ijkl}^{2}\left\vert
\psi_{i}\left(  \mathbf{\alpha}\right)  \psi_{j}\left(  \mathbf{\alpha
}\right)  \right\rangle \left\langle \psi_{k}\left(  \mathbf{\alpha}\right)
\psi_{l}\left(  \mathbf{\alpha}\right)  \right\vert
\end{equation}%
\begin{equation}
D^{2N}=\sum D_{IJ}^{2N}\left\vert \Psi_{I}\right\rangle \left\langle \Psi
_{J}\right\vert
\end{equation}%
\begin{align}
&  \mathcal{E}\left(  D^{2N},\mathbf{\alpha}\right) \\
&  =\frac{1}{\operatorname{Tr}\left\{  V\left(  \mathbf{\alpha}\right)
D^{2N}V\left(  \mathbf{\alpha}\right)  ^{\dagger}\right\}  }\operatorname{Tr}%
\left\{  \left\{  \sum_{1\leq i,k\leq2s}h_{1ik}a_{i}^{\dagger}a_{k}+\frac
{1}{4}\sum_{1\leq i,j,k,l\leq2s}\left\langle ij||kl\right\rangle
a_{i}^{\dagger}a_{j}^{\dagger}a_{l}a_{k}\right\}  V\left(  \mathbf{\alpha
}\right)  ^{\dagger}D^{2N}V\left(  \mathbf{\alpha}\right)  \right\} \\
&  \frac{1}{\operatorname{Tr}\left\{  V\left(  \mathbf{\alpha}\right)
^{\dagger}D^{2N}V\left(  \mathbf{\alpha}\right)  \right\}  }\\
&  \operatorname{Tr}\left\{  \left\{  \sum_{1\leq i,k\leq2s}h_{1ik}V\left(
\mathbf{\alpha}\right)  a_{i}^{\dagger}V\left(  \mathbf{\alpha}\right)
^{\dagger}V\left(  \mathbf{\alpha}\right)  a_{k}V\left(  \mathbf{\alpha
}\right)  ^{\dagger}+\frac{1}{4}\sum_{1\leq i,j,k,l\leq2s}\left\langle
ij||kl\right\rangle V\left(  \mathbf{\alpha}\right)  a_{i}^{\dagger}%
a_{j}^{\dagger}a_{l}a_{k}V\left(  \mathbf{\alpha}\right)  ^{\dagger}\right\}
D^{2N}\right\}
\end{align}%
\begin{align}
Va_{i}^{\dagger}V^{\dagger}  &  =\sum a_{j}^{\dagger}V_{ji}\\
Va_{i}V^{\dagger}  &  =\sum a_{j}V_{ji}^{\ast}%
\end{align}%
\begin{equation}
\left\langle Va_{k}^{\dagger}V^{\dagger}\phi|Va_{l}^{\dagger}V^{\dagger}%
\phi\right\rangle =\left\langle Va_{k}^{\dagger}\phi|Va_{l}^{\dagger}%
\phi\right\rangle =\left\langle V\psi_{k}|V\psi_{l}\right\rangle
\end{equation}%
\begin{equation}
\left\langle V\psi_{k}|V\psi_{l}\right\rangle =\sum_{k^{\prime}l^{\prime}%
}\left\langle \psi_{k^{\prime}}V_{k^{\prime}k}|\psi_{l^{\prime}}V_{l^{\prime
}l}\right\rangle =\sum_{k^{\prime}l^{\prime}}V_{k^{\prime}k}^{\ast
}\left\langle \psi_{k^{\prime}}|\psi_{l^{\prime}}\right\rangle V_{l^{\prime}l}%
\end{equation}
This can be recast as
\begin{align}
&  \mathcal{H}\left(  \mathbf{\psi}\right)  =Tr\left\{  \sum_{1\leq j,k\leq
2s}n_{j}^{\frac{1}{2}}n_{k}^{\frac{1}{2}}\left\vert \psi_{j}\right\rangle
\left\langle \psi_{k}\right\vert h_{1jk}\right. \nonumber\\
&  \qquad\qquad+\frac{1}{2}\sum_{1\leq j,k,l,m\leq2s}n_{j}^{\frac{1}{2}}%
n_{k}^{\frac{1}{2}}n_{l}^{\frac{1}{2}}n_{m}^{\frac{1}{2}}\left\{  \left\vert
\psi_{j}\right\rangle \otimes\left\vert \psi_{k}\right\rangle -\left\vert
\psi_{k}\right\rangle \otimes\left\vert \psi_{j}\right\rangle \right\}
\nonumber\\
&  \qquad\qquad\qquad\qquad\qquad\qquad\left.  \overset{}{\left\{
\left\langle \psi_{l}\right\vert \otimes\left\langle \psi_{m}\right\vert
-\left\langle \psi_{l}\right\vert \otimes\left\langle \psi_{m}\right\vert
\right\}  \left\{  h_{2jklm}-h_{2jkml}\right\}  }\right\} \nonumber\\
&  \equiv Tr\left\{  \sum_{1\leq j,k\leq2s}\left\vert \psi_{j}\right\rangle
\left\langle \psi_{k}\right\vert K_{1jk}\right. \nonumber\\
&  +\left.  \frac{1}{2}\sum_{1\leq j,k,l,m\leq2s}\left\{  \left\vert \psi
_{j}\right\rangle \otimes\left\vert \psi_{k}\right\rangle \left\langle
\psi_{l}\right\vert \otimes\left\langle \psi_{m}\right\vert -\left\vert
\psi_{j}\right\rangle \otimes\left\vert \psi_{k}\right\rangle \left\langle
\psi_{m}\right\vert \otimes\left\langle \psi_{l}\right\vert \right\}
K_{2jklm}\right\}  ,
\end{align}
where
\begin{equation}
K_{1jk}=n_{j}^{\frac{1}{2}}n_{k}^{\frac{1}{2}}h_{1jk}%
\end{equation}
and
\begin{equation}
K_{2jklm}=n_{j}^{\frac{1}{2}}n_{k}^{\frac{1}{2}}n_{l}^{\frac{1}{2}}%
n_{m}^{\frac{1}{2}}\left\{  h_{2jklm}-h_{2jkml}\right\}  . \label{defk2}%
\end{equation}
One should note that $K_{2jklm}$ is hermitian
\begin{equation}
K_{2jkml}=K_{2mljk}^{\ast}%
\end{equation}
and antisymmetric wrt interchange of the subscripts $j$ $\leftrightarrow$ $k$
and $l$ $\leftrightarrow$ $m$ i.e.
\begin{equation}
K_{2jkml}=-K_{2jklm},K_{2kjlm}=-K_{2kjlm}\quad\text{and }K_{2kjml}=K_{2jklm}.
\end{equation}
This can be seen from definition eq. $\left(  \ref{defk2}\right)  $,
\begin{align}
h_{2abcd}  &  =\int\psi_{a}\left(  \mathbf{x}_{1}\right)  ^{\ast}\psi
_{b}\left(  \mathbf{x}_{2}\right)  ^{\ast}\frac{1}{\left\Vert \mathbf{r}%
_{1}\mathbf{-r}_{2}\right\Vert }\psi_{c}\left(  \mathbf{x}_{1}\right)
\psi_{d}\left(  \mathbf{x}_{2}\right)  d\mathbf{x}_{1}d\mathbf{x}%
_{2}\nonumber\\
&  =\left[  \int\psi_{b}\left(  \mathbf{x}_{1}\right)  \psi_{a}\left(
\mathbf{x}_{2}\right)  \frac{1}{\left\Vert \mathbf{r}_{1}\mathbf{-r}%
_{2}\right\Vert }\psi_{d}\left(  \mathbf{x}_{1}\right)  ^{\ast}\psi_{c}\left(
\mathbf{x}_{2}\right)  ^{\ast}d\mathbf{x}_{1}d\mathbf{x}_{2}\right]  ^{\ast
}=h_{2dcba}^{\ast}%
\end{align}
and
\begin{align}
h_{2abcd}  &  =\int\psi_{a}\left(  \mathbf{x}_{1}\right)  ^{\ast}\psi
_{b}\left(  \mathbf{x}_{2}\right)  ^{\ast}\frac{1}{\left\Vert \mathbf{r}%
_{1}\mathbf{-r}_{2}\right\Vert }\psi_{c}\left(  \mathbf{x}_{1}\right)
\psi_{d}\left(  \mathbf{x}_{2}\right)  d\mathbf{x}_{1}d\mathbf{x}%
_{2}\nonumber\\
&  =\int\psi_{b}\left(  \mathbf{x}_{1}\right)  ^{\ast}\psi_{a}\left(
\mathbf{x}_{2}\right)  ^{\ast}\frac{1}{\left\Vert \mathbf{r}_{1}%
\mathbf{-r}_{2}\right\Vert }\psi_{d}\left(  \mathbf{x}_{1}\right)  \psi
_{c}\left(  \mathbf{x}_{2}\right)  d\mathbf{x}_{1}d\mathbf{x}_{2}=h_{2badc}.
\end{align}%
\begin{align}
\mathcal{E}\left(  \mathbf{\psi}\right)   &  =\mathcal{H}\left(  \mathbf{\psi
}\right)  =\sum_{1\leq j\leq2s}n_{j}\left\langle \psi_{j}\right\vert
h_{1}\left.  \psi_{j}\right\rangle +\frac{1}{2}\sum_{1\leq j,k\leq2s}%
n_{j}n_{k}\left\{  \left(  \left.  \psi_{j}\psi_{k}\right\vert h_{2}\psi
_{j}\psi_{k}\right)  -\left(  \left.  \psi_{j}\psi_{k}\right\vert h_{2}%
\psi_{k}\psi_{j}\right)  \right\} \nonumber\\
&  =Tr\left\{  \sum_{1\leq j\leq2s}n_{j}\left\vert \psi_{j}\right\rangle
\left\langle \psi_{j}\right\vert h_{1jj}\right. \nonumber\\
&  +\left.  \frac{1}{2}\sum_{1\leq j,k\leq2s}n_{j}n_{k}\left\{  \left\vert
\psi_{j}\right\rangle \otimes\left\vert \psi_{k}\right\rangle \left\langle
\psi_{j}\right\vert \otimes\left\langle \psi_{k}\right\vert -\left\vert
\psi_{j}\right\rangle \otimes\left\vert \psi_{k}\right\rangle \left\langle
\psi_{k}\right\vert \otimes\left\langle \psi_{j}\right\vert \right\}  \left\{
h_{2jkjk}-h_{2jkkj}\right\}  \right\}  .
\end{align}


The functional derivatives wrt $\left\{  \psi_{\mu}^{\ast};1\leq\mu
\leq2s\right\}  $ are given by%
\begin{align}
\frac{\delta\mathcal{H}\left(  \mathbf{\psi,n}\right)  }{\delta\psi_{\mu
}^{\ast}}  &  =\sum_{1\leq j\leq2s}\left\vert \psi_{j}\right\rangle K_{1j\mu
}+Tr\left\{  \frac{1}{2}\sum_{1\leq j,k,l\leq2s}\left\{  \left\vert \psi
_{j}\right\rangle \left\{  \left\vert \psi_{k}\right\rangle \left\langle
\psi_{l}\right\vert \right\}  K_{2kjjl\mu}+\left\vert \psi_{j}\right\rangle
\left\vert \psi_{k}\right\rangle \left\langle \psi_{l}\right\vert K_{2jk\mu
l}\right\}  \right\} \nonumber\\
&  =\sum_{1\leq j\leq2s}\left\vert \psi_{j}\right\rangle K_{1j\mu}+Tr\left\{
\sum_{1\leq j,k,l\leq2s}\left\vert \psi_{j}\right\rangle \left\{  \left\vert
\psi_{k}\right\rangle \left\langle \psi_{l}\right\vert \right\}  K_{2jk\mu
l}\right\}  \equiv\sum_{1\leq j\leq2s}\left\vert \psi_{j}\right\rangle
F\left(  \mathbf{\psi}\right)  _{j\mu},
\end{align}
where
\begin{equation}
F\left(  \mathbf{\psi}\right)  _{j\mu}=K_{1j\mu}+Tr\left\{  \sum_{1\leq
k,l\leq2s}\left\vert \psi_{k}\right\rangle \left\langle \psi_{l}\right\vert
K_{2jk\mu l}\right\}  =K_{1j\mu}+\sum_{1\leq k\leq2s}K_{2jk\mu k}.
\end{equation}
This defines a Fock operator by
\begin{equation}
F\left(  \mathbf{\psi}\right)  \left(  \left\vert \psi_{1}\right\rangle
,\cdots,\left\vert \psi_{2s}\right\rangle \right)  =\left(  \left\vert
\psi_{1}\right\rangle ,\cdots,\left\vert \psi_{2s}\right\rangle \right)
\mathbb{F}\left(  \mathbf{\psi}\right)  ,
\end{equation}
where $\mathbb{F}\left(  \mathbf{\psi}\right)  $ has matrix elements $\left\{
F\left(  \mathbf{\psi}\right)  _{j\mu}\right\}  $ wrt the basis $\left\{
\left\vert \psi_{j}\right\rangle ;1\leq j\leq2s\right\}  .$

The functional derivatives wrt $\left\{  \psi_{\mu};1\leq\mu\leq2s\right\}  $
are given by%
\begin{equation}
r\left\{  \sum_{1\leq j,k\leq2s}\left\vert \psi_{j}\right\rangle \left\langle
\psi_{k}\right\vert K_{1jk}+\frac{1}{2}\sum_{1\leq j,k,l,m\leq2s}\left[
\left\vert \psi_{j}\right\rangle \otimes\left\vert \psi_{k}\right\rangle
\right]  \left[  \left\langle \psi_{l}\right\vert \otimes\left\langle \psi
_{m}\right\vert \right]  K_{2jklm}\right\}
\end{equation}%
\begin{align}
\frac{\delta\mathcal{H}\left(  \mathbf{\psi,n}\right)  }{\delta\psi_{\mu}}  &
=\sum_{1\leq j\leq2s}K_{1\mu j}\left\langle \psi_{\mu}\right\vert +Tr\left\{
\frac{1}{2}\sum_{1\leq j,l,m\leq2s}\left\vert \psi_{j}\right\rangle \left[
\left\langle \psi_{l}\right\vert \otimes\left\langle \psi_{m}\right\vert
\right]  K_{2\mu jlm}\right\} \nonumber\\
&  \sum_{1\leq j,l,m\leq2s}\left[  \left\langle \psi_{l}\right\vert
\otimes\left[  \left\vert \psi_{j}\right\rangle \right]  \left\langle \psi
_{m}\right\vert \right]  K_{2j\mu lm}\\
&  =\left\{  K_{1\mu\mu}+Tr\left\{  \sum_{1\leq k\leq2s}K_{2\mu k\mu
k}\right\}  \right\}  \left\langle \psi_{\mu}\right\vert \equiv\left\langle
\psi_{\mu}\right\vert F\left(  \mathbf{\psi}\right)  _{\mu\mu}.
\end{align}%
\begin{align}
a_{p}^{\dagger}\sum_{1\leq i,k\leq2s}h_{1ik}\left[  a_{q},a_{i}^{\dagger}%
a_{k}\right]   &  =a_{p}^{\dagger}\sum_{1\leq i,k\leq2s}h_{1ik}\left(
a_{q}a_{i}^{\dagger}a_{k}-a_{i}^{\dagger}a_{k}a_{q}\right) \nonumber\\
&  =a_{p}^{\dagger}\sum_{1\leq i,k\leq2s}h_{1ik}\left(  \left(  \delta
_{qi}-a_{i}^{\dagger}a_{q}\right)  a_{k}-a_{i}^{\dagger}a_{k}a_{q}\right)
\nonumber\\
&  =a_{p}^{\dagger}\sum_{1\leq i,k\leq2s}h_{1ik}a_{k}\delta_{qi}%
=a_{p}^{\dagger}\sum_{1\leq k\leq2s}h_{1qk}a_{k}=\sum_{1\leq k\leq2s}%
h_{1qk}a_{p}^{\dagger}a_{k}%
\end{align}%
\begin{align}
&  a_{p}^{\dagger}\sum_{1\leq i,j,k,l\leq2s}V_{2ijkl}\left[  a_{q}%
,a_{i}^{\dagger}a_{j}^{\dagger}a_{l}a_{k}\right]  =a_{p}^{\dagger}\sum_{1\leq
i,j,k,l\leq2s}V_{2ijkl}\left(  a_{q}a_{i}^{\dagger}a_{j}^{\dagger}a_{l}%
a_{k}-a_{i}^{\dagger}a_{j}^{\dagger}a_{l}a_{k}a_{q}\right) \nonumber\\
&  =a_{p}^{\dagger}\sum_{1\leq i,j,k,l\leq2s}V_{2ijkl}\left(  \left(
\delta_{qi}-a_{i}^{\dagger}a_{q}\right)  a_{j}^{\dagger}a_{l}a_{k}%
-a_{i}^{\dagger}a_{j}^{\dagger}a_{q}a_{l}a_{k}\right) \nonumber\\
&  =a_{p}^{\dagger}\sum_{1\leq i,j,k,l\leq2s}V_{2ijkl}\left(  \left(
\delta_{qi}a_{j}^{\dagger}a_{l}a_{k}-a_{i}^{\dagger}\left(  \delta_{qj}%
-a_{j}^{\dagger}a_{q}\right)  a_{l}a_{k}\right)  -a_{i}^{\dagger}%
a_{j}^{\dagger}a_{q}a_{l}a_{k}\right) \nonumber\\
&  =a_{p}^{\dagger}\sum_{1\leq i,j,k,l\leq2s}V_{2ijkl}\left(  \delta_{qi}%
a_{j}^{\dagger}a_{l}a_{k}-a_{i}^{\dagger}a_{l}a_{k}\delta_{qj}+a_{i}^{\dagger
}a_{j}^{\dagger}a_{q}a_{l}a_{k}-a_{i}^{\dagger}a_{j}^{\dagger}a_{q}a_{l}%
a_{k}\right) \nonumber\\
&  =a_{p}^{\dagger}\sum_{1\leq i,j,k,l\leq2s}V_{2ijkl}\left(  \delta_{qi}%
a_{j}^{\dagger}a_{l}a_{k}-a_{i}^{\dagger}a_{l}a_{k}\delta_{qj}\right)
=a_{p}^{\dagger}\sum_{1\leq i,j,k,l\leq2s}\left(  V_{2qjkl}a_{j}^{\dagger
}a_{l}a_{k}-V_{2jqkl}a_{j}^{\dagger}a_{l}a_{k}\right) \nonumber\\
&  =2\sum_{1\leq j,k,l\leq2s}V_{2qjkl}a_{p}^{\dagger}a_{j}^{\dagger}a_{l}a_{k}%
\end{align}
The AGP\ energy functional can be written as
\begin{align}
\mathcal{E}\left(  \mathbf{\eta,\alpha}\right)   &  =\frac{1}%
{\operatorname{Tr}\left\{  V\left(  \mathbf{\eta}\right)  ^{\dagger}D_{0}%
^{2N}V\left(  \mathbf{\eta}\right)  \right\}  }\operatorname{Tr}\left\{
\left\{  \sum_{1\leq i,k\leq2s}\left\langle \left.  \varphi_{i}\left(
\mathbf{\alpha}\right)  \right\vert h_{1}\varphi_{k}\left(  \mathbf{\alpha
}\right)  \right\rangle a_{i}^{\dagger}a_{k}\right.  \right. \nonumber\\
&  \qquad\qquad\qquad+\left.  \left.  \frac{1}{4}\sum_{1\leq i,j,k,l\leq
2s}\left\langle \left.  \varphi_{i}\left(  \mathbf{\alpha}\right)  \varphi
_{j}\left(  \mathbf{\alpha}\right)  \right\vert h_{2}\varphi_{k}\left(
\mathbf{\alpha}\right)  \varphi_{l}\left(  \mathbf{\alpha}\right)
\right\rangle a_{i}^{\dagger}a_{j}^{\dagger}a_{l}a_{k}\right\}  V\left(
\mathbf{\eta}\right)  D_{0}^{2N}V\left(  \mathbf{\eta}\right)  \right\}  ,
\label{agpen1}%
\end{align}
where $V\left(  \mathbf{\eta}\right)  $ are the self adjoint positive
transformations
\begin{equation}
V\left(  \mathbf{\eta}\right)  =\exp\left\{  \sum_{1\leq j\leq s}\eta
_{j}\left(  a_{j}^{\dagger}a_{j}+a_{j+s}^{\dagger}a_{j+s}\right)  \right\}
;\quad\eta_{j}\in\mathbb{R}%
\end{equation}
and the unitary group parameters $\mathbf{\alpha}$ include $\left\{
\mathbf{\lambda},\mathbf{\theta}\right\}  ,$ as well as the parameters
$\mathbf{\kappa}$ that generate the invariance group eq. $\left(
\ref{iso}\right)  .$ Unlike the general case one can explicitly parameterize
$N$-representable AGP\ SORDM's with matrix elements
\begin{equation}
\mathbb{D}^{2}\left(  \mathbf{\eta}\right)  \leftrightarrow\left\{  \frac
{1}{\operatorname{Tr}\left\{  V\left(  \mathbf{\eta}\right)  D_{0}%
^{2N}V\left(  \mathbf{\eta}\right)  \right\}  }\operatorname{Tr}\left\{
a_{i}^{\dagger}a_{j}^{\dagger}a_{l}a_{k}V\left(  \mathbf{\eta}\right)
D_{0}^{2N}V\left(  \mathbf{\eta}\right)  \right\}  \right\}
\end{equation}
that are given in terms of the geminal occupation numbers $\mathbf{n}%
=\mathbf{n}\left(  \mathbf{\eta}\right)  .$ The energy expression eq. $\left(
\ref{agpen1}\right)  $ can be manipulated to give%
\begin{align}
&  \mathcal{E}\left(  \mathbf{\eta,\alpha}\right)  =\frac{1}{\operatorname{Tr}%
\left\{  V\left(  \mathbf{\eta}\right)  D_{0}^{2N}V\left(  \mathbf{\eta
}\right)  \right\}  }\operatorname{Tr}\left\{  \left\{  \sum_{1\leq i,k\leq
2s}\left\langle \left.  \varphi_{i}\left(  \mathbf{\alpha}\right)  \right\vert
h_{1}\varphi_{k}\left(  \mathbf{\alpha}\right)  \right\rangle V\left(
\mathbf{\eta}\right)  a_{i}^{\dagger}V\left(  \mathbf{\eta}\right)
^{-1}V\left(  \mathbf{\eta}\right)  a_{k}V\left(  \mathbf{\eta}\right)
\right.  \right. \nonumber\\
&  \qquad\qquad\qquad\qquad\qquad+\frac{1}{4}\sum_{1\leq i,j,k,l\leq
2s}\left\langle \left.  \varphi_{i}\left(  \mathbf{\alpha}\right)  \varphi
_{j}\left(  \mathbf{\alpha}\right)  \right\vert h_{2}\varphi_{k}\left(
\mathbf{\alpha}\right)  \varphi_{l}\left(  \mathbf{\alpha}\right)
\right\rangle \times\\
&  \qquad\qquad\qquad\qquad\qquad\left.  \left.  V\left(  \mathbf{\eta
}\right)  a_{i}^{\dagger}V\left(  \mathbf{\eta}\right)  ^{-1}V\left(
\mathbf{\eta}\right)  a_{j}^{\dagger}V\left(  \mathbf{\eta}\right)
^{-1}V\left(  \mathbf{\eta}\right)  a_{l}V\left(  \mathbf{\eta}\right)
V\left(  \mathbf{\eta}\right)  ^{-1}a_{k}V\left(  \mathbf{\eta}\right)
\right\}  D_{0}^{2N}\right\}
\end{align}%
\begin{align}
\mathcal{E}\left(  \mathbf{\eta,\alpha}\right)   &  =\frac{1}%
{\operatorname{Tr}\left\{  V\left(  \mathbf{\eta}\right)  D_{0}^{2N}V\left(
\mathbf{\eta}\right)  \right\}  }\operatorname{Tr}\left\{  \left\{
\sum_{1\leq i,k\leq2s}\left\langle \left.  \varphi_{i}\left(  \mathbf{\alpha
}\right)  \right\vert h_{1}\varphi_{k}\left(  \mathbf{\alpha}\right)
\right\rangle e^{\eta_{i}\widehat{n_{i}}}a_{i}^{\dagger}a_{k}e^{\eta
_{k}\widehat{n_{k}}}\right.  \right. \nonumber\\
&  \qquad+\left.  \left.  \frac{1}{4}\sum_{1\leq i,j,k,l\leq2s}\left\langle
\left.  \varphi_{i}\left(  \mathbf{\alpha}\right)  \varphi_{j}\left(
\mathbf{\alpha}\right)  \right\vert h_{2}\varphi_{k}\left(  \mathbf{\alpha
}\right)  \varphi_{l}\left(  \mathbf{\alpha}\right)  \right\rangle e^{\eta
_{i}\widehat{n_{i}}}a_{i}^{\dagger}e^{\eta_{j}\widehat{n_{j}}}a_{j}^{\dagger
}a_{l}e^{\eta_{l}\widehat{n_{l}}}a_{k}e^{\eta_{k}\widehat{n_{k}}}\right\}
D_{0}^{2N}\right\}  ,
\end{align}%
\begin{align}
\mathcal{E}\left(  \mathbf{\eta,\alpha}\right)   &  =\frac{1}%
{\operatorname{Tr}\left\{  V\left(  \mathbf{\eta}\right)  D_{0}^{2N}V\left(
\mathbf{\eta}\right)  \right\}  }\operatorname{Tr}\left\{  \left\{
\sum_{1\leq i,k\leq2s}\left\langle \left.  \varphi_{i}\left(  \mathbf{\alpha
}\right)  \right\vert h_{1}\varphi_{k}\left(  \mathbf{\alpha}\right)
\right\rangle e^{\eta_{i}}a_{i}^{\dagger}a_{k}e^{\eta_{k}}\right.  \right.
\nonumber\\
&  \qquad+\left.  \left.  \frac{1}{4}\sum_{1\leq i,j,k,l\leq2s}\left\langle
\left.  \varphi_{i}\left(  \mathbf{\alpha}\right)  \varphi_{j}\left(
\mathbf{\alpha}\right)  \right\vert h_{2}\varphi_{k}\left(  \mathbf{\alpha
}\right)  \varphi_{l}\left(  \mathbf{\alpha}\right)  \right\rangle e^{\eta
_{i}}a_{i}^{\dagger}e^{\eta_{j}}a_{j}^{\dagger}a_{l}e^{\eta_{l}}a_{k}%
e^{\eta_{k}}\right\}  D_{0}^{2N}\right\}  ,
\end{align}%
\begin{align}
\mathcal{E}\left(  \mathbf{\eta,\alpha}\right)   &  =\frac{1}%
{\operatorname{Tr}\left\{  V\left(  \mathbf{\eta}\right)  D_{0}^{2N}V\left(
\mathbf{\eta}\right)  \right\}  }\operatorname{Tr}\left\{  \left\{
\sum_{1\leq i,k\leq2s}\left\langle \left.  e^{\eta_{i}}\varphi_{i}\left(
\mathbf{\alpha}\right)  \right\vert h_{1}e^{\eta_{k}}\varphi_{k}\left(
\mathbf{\alpha}\right)  \right\rangle a_{i}^{\dagger}a_{k}\right.  \right.
\nonumber\\
&  \qquad+\left.  \left.  \frac{1}{4}\sum_{1\leq i,j,k,l\leq2s}\left\langle
\left.  e^{\eta_{i}}\varphi_{i}\left(  \mathbf{\alpha}\right)  e^{\eta_{j}%
}\varphi_{j}\left(  \mathbf{\alpha}\right)  \right\vert h_{2}e^{\eta_{l}%
}\varphi_{k}\left(  \mathbf{\alpha}\right)  e^{\eta_{k}}\varphi_{l}\left(
\mathbf{\alpha}\right)  \right\rangle a_{i}^{\dagger}a_{j}^{\dagger}a_{l}%
a_{k}\right\}  D_{0}^{2N}\right\}  ,
\end{align}%
\begin{align}
\mathcal{E}\left(  \mathbf{\chi}\left(  \mathbf{\eta,\alpha}\right)  \right)
&  =\frac{1}{\operatorname{Tr}\left\{  D^{2N}\left(  \mathbf{\chi}\left(
\mathbf{\eta,\alpha}\right)  \right)  \right\}  }\operatorname{Tr}\left\{
\left\{  \sum_{1\leq i,k\leq2s}\left\langle \left.  \chi_{i}\left(
\mathbf{\eta,\alpha}\right)  \right\vert h_{1}\chi_{k}\left(  \mathbf{\eta
,\alpha}\right)  \right\rangle a_{i}^{\dagger}a_{k}\right.  \right.
\nonumber\\
&  \qquad+\left.  \left.  \frac{1}{4}\sum_{1\leq i,j,k,l\leq2s}\left\langle
\left.  \chi_{i}\left(  \mathbf{\eta,\alpha}\right)  \chi_{j}\left(
\mathbf{\eta,\alpha}\right)  \right\vert h_{2}\chi_{k}\left(  \mathbf{\eta
,\alpha}\right)  \chi_{l}\left(  \mathbf{\eta,\alpha}\right)  \right\rangle
a_{i}^{\dagger}a_{j}^{\dagger}a_{l}a_{k}\right\}  D_{0}^{2N}\right\}  ,
\end{align}
We can either vary the parameters $\left(  \mathbf{\eta,\alpha}\right)  $ in
an unconstrained fashion or vary the GSO's $\left\{  \left.  \chi
_{j}\right\vert 1\leq j\leq2s\right\}  $ in the energy functional
\begin{align}
\mathcal{E}\left(  \mathbf{\chi}\right)   &  =\frac{1}{\operatorname{Tr}%
\left\{  D^{2N}\left(  \mathbf{\chi}\right)  \right\}  }\operatorname{Tr}%
\left\{  \left\{  \sum_{1\leq i,k\leq2s}\left\langle \left.  \chi
_{i}\right\vert h_{1}\chi_{k}\right\rangle a_{i}^{\dagger}a_{k}\right.
\right. \nonumber\\
&  \qquad+\left.  \left.  \frac{1}{4}\sum_{1\leq i,j,k,l\leq2s}\left\langle
\left.  \chi_{i}\chi_{j}\right\vert h_{2}\chi_{k}\chi_{l}\right\rangle
a_{i}^{\dagger}a_{j}^{\dagger}a_{l}a_{k}\right\}  D_{0}^{2N}\right\}  ,
\end{align}
o that they are constrained to be orthogonal, but \emph{not} normalized. In
order to do this we consider the Lagrangian
\begin{equation}
\mathcal{L}\left(  \mathbf{\chi}\right)  =\mathcal{E}\left(  \mathbf{\chi
}\right)  -\sum_{j\neq k}\varepsilon_{kj}\left\langle \left.  \chi
_{j}\right\vert \chi_{k}\right\rangle ,
\end{equation}
which gives the stationarity condition
\begin{equation}
\frac{\delta\mathcal{L}\left(  \mathbf{\chi}\right)  }{\delta\chi_{\mu}^{\ast
}}=\frac{\delta\mathcal{E}\left(  \mathbf{\chi}\right)  }{\delta\chi_{\mu
}^{\ast}}-\sum_{j\neq\mu}\varepsilon_{j\mu}\left\vert \chi_{j}\right\rangle
=\frac{1}{S_{N}\left(  \mathbf{\chi}\right)  }\left\{  \frac{\delta
\mathcal{H}\left(  \mathbf{\chi}\right)  }{\delta\chi_{\mu}^{\ast}%
}-\mathcal{E}\left(  \mathbf{\chi}\right)  \frac{\delta S_{N}\left(
\mathbf{\chi}\right)  }{\delta\chi_{\mu}^{\ast}}\right\}  -\sum_{j\neq\mu
}\varepsilon_{j\mu}\left\vert \psi_{j}\right\rangle =0.
\end{equation}%
\begin{align}
F\left(  \mathbf{\psi}\right)   &  =\sum_{1\leq m,n\leq2s}\left\{  \left(
\mathbb{H}_{1}\left(  \mathbf{\psi}\right)  \mathbb{D}_{0}^{1}\right)
_{mn}+\frac{1}{2}\left(  L_{2}^{1}\left(  \mathbb{V}_{2}\left(  \mathbf{\psi
}\right)  \mathbb{D}_{0}^{2}\right)  \right)  _{mn}-\mathcal{E}\left(
\mathbf{\psi}\right)  D_{0mm}^{1}\delta_{mn}\right\}  \frac{\left\vert
\psi_{m}\right\rangle \left\langle \psi_{n}\right\vert }{\left\Vert \psi
_{m}\right\Vert \left\Vert \psi_{n}\right\Vert },\\
&  \Rightarrow F\left(  \mathbf{\psi}\right)  _{mm}=\left(  \mathbb{H}%
_{1}\left(  \mathbf{\psi}\right)  \mathbb{D}_{0}^{1}\right)  _{mm}+\frac{1}%
{2}\left(  L_{2}^{1}\left(  \mathbb{V}_{2}\left(  \mathbf{\psi}\right)
\mathbb{D}_{0}^{2}\right)  \right)  _{mm}-\mathcal{E}\left(  \mathbf{\psi
}\right)  D_{0mm}^{1}=0\\
F\left(  \mathbf{\psi}\right)  _{mn}  &  =\left(  \mathbb{H}_{1}\left(
\mathbf{\psi}\right)  \mathbb{D}_{0}^{1}\right)  _{mn}+\frac{1}{2}\left(
L_{2}^{1}\left(  \mathbb{V}_{2}\left(  \mathbf{\psi}\right)  \mathbb{D}%
_{0}^{2}\right)  \right)  _{mn}=\mathbb{H}_{1}\left(  \mathbf{\psi}\right)
_{mn}\mathbb{D}_{0nn}^{1}+\frac{1}{2}\left(  L_{2}^{1}\left(  \mathbb{V}%
_{2}\left(  \mathbf{\psi}\right)  \mathbb{D}_{0}^{2}\right)  \right)  _{mn}%
\end{align}%
\begin{equation}
\frac{\delta\mathcal{E}\left(  \mathbf{\psi}\right)  }{\delta\psi_{\mu}^{\ast
}}=F\left(  \mathbf{\psi}\right)  \left\vert \psi_{\mu}\right\rangle
\end{equation}%
\begin{equation}
\left\langle \left.  \Psi\right\vert Ha_{i}^{\dagger}a_{j}\Psi\right\rangle
=0;\quad1\leq i\neq j\pm s\leq2s
\end{equation}
as
\begin{align}
&  \left\langle g_{0}^{N}\right.  \left\vert \left[  H,q_{ij}^{\dagger
}\right]  g_{0}^{N}\right\rangle =0;\quad1\leq i\neq j\pm s\leq2s\\
&  \Rightarrow\left\langle g_{0}^{N}\right.  \left\vert Hq_{ij}^{\dagger}%
g_{0}^{N}\right\rangle =0\\
&  q_{ij}\left\vert g_{0}^{N}\right\rangle =0\\
&  \left\langle g_{0}^{N}\right.  \left\vert H\sum_{1\leq i\neq j\pm s\leq
2s}\left\{  \zeta_{ij}q_{ij}^{\dagger}+\zeta_{ij}^{\ast}q_{ij}\right\}
g_{0}^{N}\right\rangle =0;
\end{align}%
\begin{align}
\left\langle \left.  \Psi\right\vert \left[  H,a_{i}^{\dagger}a_{j}\right]
\Psi\right\rangle  &  =0;\quad1\leq i=j\pm s\leq2s\\
\left\langle \left.  \Psi\right\vert Ha_{i}^{\dagger}a_{i\pm s}\Psi
\right\rangle  &  =0
\end{align}%
\begin{equation}
\left\langle g_{0}^{N}\right.  \left\vert H\hat{n}_{1j}g_{0}^{N}\right\rangle
=\mathcal{E}\left\langle g_{0}^{N}\right.  \left\vert \hat{n}_{j1}g_{0}%
^{N}\right\rangle
\end{equation}
gives the stationary occupation number condition. This comes from
\begin{equation}
\frac{1}{S_{N}}\left\{  \left\langle g_{0}^{N}\right.  \left\vert \left(
\hat{n}_{1j}H+H\hat{n}_{1j}\right)  g_{0}^{N}\right\rangle -2\mathcal{E}%
\left\langle g_{0}^{N}\right.  \left\vert \hat{n}_{1j}g_{0}^{N}\right\rangle
\right\}  =0.
\end{equation}
where $\left\{  \left\vert \psi_{j}\right\rangle \right\}  $ are \emph{any
}one electron orthogonal states normalized or unnormalized belonging to
$\mathcal{H}^{1}.$The variational procedure can be schematically displayed as
\begin{align}
&  \left(  \mathbf{n}_{i},\mathbf{\theta}_{i}\mathbf{,}\left\vert
\mathbf{\varphi}_{i}\right\rangle \right)  \overset{\text{ }}{\qquad
\qquad\qquad\underrightarrow{\text{iterate on the Fock Operator}}\qquad\left(
\mathbf{n}_{i},\mathbf{\theta}_{i}\mathbf{,}\left\vert \mathbf{\varphi}%
_{i+1}\left(  \mathbf{n}_{i},\mathbf{\theta}_{i}\right)  \right\rangle
\right)  }\nonumber\\
&  \qquad\left\uparrow
\begin{array}
[c]{c}%
\\
i\rightarrow i+1\\
\end{array}
\right.  \qquad\qquad\qquad\qquad\qquad\qquad\qquad\qquad\qquad
\left\downarrow
\begin{array}
[c]{c}%
\\
\text{minimize wrt }\mathbf{\theta}\\
\end{array}
\right. \nonumber\\
&  \left(  \mathbf{n}_{i+1},\mathbf{\theta}_{i+1}\mathbf{,}\left\vert
\mathbf{\varphi}_{i+1}\left(  \mathbf{n}_{i},\mathbf{\theta}_{i+1}\right)
\right\rangle \right)  \underleftarrow{\qquad\text{minimize wrt }%
\mathbf{n}\qquad\qquad}\left(  \mathbf{n}_{i},\mathbf{\theta}_{i+1}%
\mathbf{,}\left\vert \mathbf{\varphi}_{i+1}\left(  \mathbf{n}_{i}%
,\mathbf{\theta}_{i+1}\right)  \right\rangle \right)  .
\end{align}
The GSO's and $\mathbf{\theta}$ that are self consistent solutions to the
above scheme determine the CGSO's and together with the self consistent
occupation numbers $\mathbf{n}$ determine the optimal geminal and
$2N$-electron AGP state.

\subsection{Orthonormality Constraints}

A necessary condition for $\mathbf{\psi}_{\#}$ $=\mathbf{\psi}_{\#}\left(
\mathbf{n,\theta}\right)  $ to be a orthonormally constrained stationary point
of $\mathcal{E}\left(  \mathbf{n,\psi,\theta}\right)  $ for a fixed value of
$\left(  \mathbf{n,\theta}\right)  $ is given by
\begin{subequations}
\begin{align}
\partial_{\mathbf{\psi}}^{1}\mathcal{\tilde{H}}\left(  \mathbf{n,\mathbf{\psi
}_{\#},\theta}\right)   &  =\mathbf{0}\label{nec1}\\
\left\langle \left.  \psi_{\#j}\right\vert \psi_{\#k}\right\rangle  &
=\delta_{jk};1\leq j,k\leq2s,
\end{align}
where the Lagrangian $\mathcal{\tilde{H}}\left(  \mathbf{n,\mathbf{\psi
},\theta}\right)  $ is defined as
\end{subequations}
\begin{equation}
\mathcal{\tilde{H}}\left(  \mathbf{n,\mathbf{\psi},\theta}\right)
=\mathcal{H}\left(  \mathbf{n,\mathbf{\psi},\theta}\right)  -\sum_{1\leq
i,j\leq2s}\varepsilon_{ij}\left\{  \left\langle \left.  \psi_{i}\right\vert
\psi_{j}\right\rangle -1\right\}  , \label{lagH}%
\end{equation}
as
\begin{equation}
\partial_{\mathbf{\psi}}^{1}\mathcal{\tilde{H}}\left(  \mathbf{n,\mathbf{\psi
},\theta}\right)  =\partial_{\mathbf{\psi}}^{1}\mathcal{\tilde{E}}\left(
\mathbf{n,\mathbf{\psi},\theta}\right)
\end{equation}
for orthonormally constrained $\mathbf{\mathbf{\psi.}}$ The condition eq.
$\left(  \ref{nec1}\right)  $ can be expressed in terms of the generalized
Fock operator as
\begin{equation}
F\left(  \mathbf{n,\mathbf{\psi},\theta}\right)  \left\vert \psi
_{j}\right\rangle =\varepsilon_{j}\left\vert \psi_{j}\right\rangle ;1\leq
j\leq2s. \label{eig}%
\end{equation}
If $F\left(  \mathbf{n,\mathbf{\psi},\theta}\right)  =F\left(
\mathbf{n,\mathbf{\psi},\theta}\right)  ^{\dagger}$, (i.e. the state defined
by $\left(  \mathbf{n,\mathbf{\psi},\theta}\right)  $ is stationary), this
implies that one can find solutions $\left\{  \left.  \left\vert \psi
_{m}\right\rangle \right\vert 1\leq m\leq2s\right\}  $ for $\left\{  \left.
\varepsilon_{m}\right\vert 1\leq m\leq2s\right\}  $ such that
\begin{equation}
\left\langle \left.  \psi_{j}\right\vert \psi_{k}\right\rangle =\delta
_{jk},;\quad1\leq j,k\leq2s.
\end{equation}
Thus one can always find stationary GSO's that satisfy the \emph{normality}
constraints
\begin{equation}
\left\langle \left.  \psi_{j}\right\vert \psi_{j}\right\rangle =1;\quad1\leq
j\leq2s
\end{equation}
that \emph{also} satisfy the orthonormality constraints eq. $\left(
\ref{norm}\right)  .$ This conclusion is consistent with the observation that
the orthogonality constraints are \emph{not }active for $\left\{  \left.
\left\vert \psi_{j}\right\rangle \right\vert \right\}  $ that satisfy eq.
$\left(  \ref{eig}\right)  $ as "off diagonal" Lagrange multipliers,$\left\{
\left.  \varepsilon_{jk}\right\vert j\neq k\right\}  $ that are associated
with such constraints have the value zero.%

\begin{align}
F\left(  \mathbf{n},\mathbf{\psi}\right)  _{pp} &  =\frac{1}{S_{N}\left(
g^{N}\right)  }\left\langle \left.  g\right\vert a_{p}^{\dagger}a_{p}%
H-a_{p}^{\dagger}Ha_{p}g\right\rangle =\mathcal{E}\left(  g\right)
N_{p}-\frac{1}{S_{N}\left(  g\right)  }\left\langle \left.  a_{p}g\right\vert
Ha_{p}g\right\rangle \\
&  =\mathcal{E}\left(  g\right)  N_{p}-N_{p}\left\langle \left.  \psi
_{p}\right\vert H\psi_{p}\right\rangle =N_{p}\left(  \mathcal{E}\left(
g\right)  -\left\langle \left.  \psi_{p}\right\vert H\psi_{p}\right\rangle
\right)  \\
&  \left(  \sum\left\langle \left.  \psi_{q}\right\vert H\psi_{q}\right\rangle
N_{q}+\mathcal{H}_{2}-\left\langle \left.  \psi_{p}\right\vert H\psi
_{p}\right\rangle \right)  =\left(  N_{p}-1\right)  \left\langle \left.
\psi_{p}\right\vert H\psi_{p}\right\rangle +\sum_{q\neq p}\left\langle \left.
\psi_{q}\right\vert H\psi_{q}\right\rangle N_{q}+\mathcal{H}_{2}%
\end{align}%
\begin{align}
&  \frac{\delta\mathcal{E}\left(  \mathbf{n,\psi}\right)  }{\delta\psi
_{j}^{\ast}}=F\left(  \mathbf{n,\psi}\right)  \left\vert \psi_{j}\right\rangle
=\sum_{p}\left\vert \psi_{p}\right\rangle F\left(  \mathbf{n},\mathbf{\psi
}\right)  _{pj}\equiv\sum\left\vert \psi_{p}\right\rangle \varepsilon
_{pj}\Rightarrow\\
&  \left\vert \psi_{j}\right\rangle N_{j}\left\{  \mathcal{E}\left(
g^{N}\right)  -\mathcal{E}_{2N-1}\left(  \psi_{j}\right)  \right\}
-\sum_{p\neq j}\left\vert \psi_{p}\right\rangle \frac{1}{S_{N}\left(
g^{N}\right)  }\left\langle \left.  a_{p}g^{N}\right\vert Ha_{j}%
g^{N}\right\rangle
\end{align}
In Appendix \ref{saddle} we show that a Fock matrix associated with a state
minimized by orthonormally constrained orbitals is negative, which in
particular leads to $F\left(  \mathbf{n},\mathbf{\psi}\right)  _{pp}\leq0$ for
$1\leq p\leq2s$. Hence
\begin{equation}
I\left(  \psi_{p}\right)  =\mathcal{E}\left(  g^{N}\right)  -\mathcal{E}%
_{2N-1}\left(  \psi_{p}\right)  <0,
\end{equation}
where $I\left(  \psi_{p}\right)  $ is the ionization potential associated with
the CGSO $\left\vert \psi_{p}\right\rangle $.In Appendix \ref{aufbau} we
deduce that in order to minimize $\mathcal{E}\left(  \mathbf{n},\mathbf{\psi
}\right)  $ the highest occupation numbers must be associated with the most
negative values of $\varsigma\left(  \psi_{l}\right)  ,$ hence leading to a
generalized aufbau principle.%

\begin{equation}
=F\left(  g\left(  \mathbf{\eta},\mathbf{\psi}\right)  ^{N}\right)
_{q+sq}=F\left(  g\left(  \mathbf{\eta},\mathbf{\psi}\right)  ^{N}\right)
_{qq+s}%
\end{equation}
The normalization $\operatorname{Tr}\left\{  D^{2N}\right\}  $ does not depend
on $\left\{  \varphi_{j}\right\}  $ as they are constrained to be orthonormal
thus%
\begin{equation}
\frac{\delta\mathcal{E}\left(  \mathbb{D}^{2},\mathbf{\varphi}\right)
}{\delta\varphi_{j}^{\ast}\left(  \mathbf{r,\xi}\right)  }=\frac
{1}{\operatorname{Tr}\left\{  D^{2N}\right\}  }\frac{\delta\mathcal{H}\left(
\mathbb{D}^{2},\mathbf{\varphi}\right)  }{\delta\varphi_{j}^{\ast}\left(
\mathbf{r,\xi}\right)  }%
\end{equation}%
\begin{equation}
\mathcal{E}\left(  \mathbb{D}^{2}\left(  \eta\right)  ,\mathbf{\varphi
}\right)  =\frac{1}{2\left(  N-1\right)  }\left\{  \sum_{1\leq p\leq2s}%
\frac{\operatorname{Tr}\left\{  a_{p}^{\dagger}a_{p}k\left(  \eta\right)
D^{2N}k\left(  \eta\right)  \right\}  }{\operatorname{Tr}\left\{  k\left(
\eta\right)  D^{2N}k\left(  \eta\right)  \right\}  }\left\{  \frac
{\operatorname{Tr}\left\{  a_{p}^{\dagger}Ha_{p}k\left(  \eta\right)
D^{2N}k\left(  \eta\right)  \right\}  }{\operatorname{Tr}\left\{
a_{p}^{\dagger}a_{p}k\left(  \eta\right)  D^{2N}k\left(  \eta\right)
\right\}  }-\left\langle \left.  \varphi_{p}\right\vert h_{1}\varphi
_{p}\right\rangle \right\}  \right\}
\end{equation}%
\begin{equation}
\mathcal{E}\left(  \mathbb{D}^{2},\mathbf{\varphi}\right)  =\frac{1}{2\left(
N-1\right)  }\left\{  \sum_{1\leq p\leq2s}N_{p}\left(  \mathbb{D}%
^{2},\mathbf{\varphi}\right)  \left\{  \frac{\operatorname{Tr}\left\{
a_{p}^{\dagger}Ha_{p}D^{2N}\right\}  }{\operatorname{Tr}\left\{
a_{p}^{\dagger}a_{p}D^{2N}\right\}  }-\left\langle \left.  \varphi
_{p}\right\vert h_{1}\varphi_{p}\right\rangle \right\}  \right\}
\end{equation}
Aufbau Principle\label{aufbau}

The energy of any AGP state can be expressed in terms of weighted ionization
potentials $\left\{  N_{p}\left(  \mathbf{\psi}\right)  I\left(  \psi
_{p}\right)  \right\}  $ and the one electron energy $\mathcal{E}_{1}\left(
\mathbf{n},\mathbf{\psi}\right)  $ as
\begin{equation}
\mathcal{E}\left(  \mathbf{n},\mathbf{\psi}\right)  =\frac{1}{2}\left\{
\sum_{1\leq p\leq2s}N_{p}\left(  \mathbf{\psi}\right)  I\left(  \psi
_{p}\right)  +\mathcal{E}_{1}\left(  \mathbf{n},\mathbf{\psi}\right)
\right\}  .
\end{equation}
Let $\left\vert g^{N}\left(  \mathbf{n},\mathbf{\psi}\right)  \right\rangle $
be a \emph{globally minimized }AGP state such that \emph{\ }
\begin{equation}
I\left(  \psi_{a}\right)  <I\left(  \psi_{b}\right)
\end{equation}
and
\begin{equation}
N_{b}>N_{a}.
\end{equation}
We will compare the energy $\mathcal{E}\left(  \mathbf{n},\mathbf{\psi
}\right)  $ with the energy $\mathcal{E}\left(  \mathbf{n}_{ba},\mathbf{\psi
}\right)  $ produced when the occupation numbers $n_{a}$ and $n_{b}$ are
interchanged, but everything else is left constant. Interchanging $n_{a}$ with
$n_{b}$ interchanges $N_{a}$ and $N_{b}$ as $\mathbf{n}$ is monotonically in
$1-1$ correspondence with $N$. Changing the occupations numbers
\begin{equation}
\left(  n_{a}\leftrightarrow n_{b}\Leftrightarrow N_{a}\leftrightarrow
N_{b}\right)
\end{equation}
can be produced by an occupation number transformation
\begin{equation}
\left\vert g^{N}\left(  \mathbf{n}_{ba},\mathbf{\psi}\right)  \right\rangle
=k\left(  \mathbf{\eta}_{ba}\right)  \left\vert g^{N}\left(  \mathbf{n}%
,\mathbf{\psi}\right)  \right\rangle .
\end{equation}
to produce the new energy
\begin{equation}
\mathcal{E}\left(  \mathbf{n}_{ba},\mathbf{\psi}\right)  =\frac{1}{2}%
%TCIMACRO{\tsum \limits_{1\leq p\leq2s}}%
%BeginExpansion
{\textstyle\sum\limits_{1\leq p\leq2s}}
%EndExpansion
\left\{  \frac{1}{S_{N}\left(  \mathbf{n}_{ba},\mathbf{\psi}\right)
}\left\langle \left.  g^{N}\left(  \mathbf{n}_{ba},\mathbf{\psi}\right)
\right\vert \left\{  a_{p}^{\dagger}a_{p}H-a_{p}^{\dagger}Ha_{p}\right\}
g^{N}\left(  \mathbf{n}_{ba},\mathbf{\psi}\right)  \right\rangle
+N_{p}\left\langle \left.  \psi_{p}\right\vert h_{1}\psi_{p}\right\rangle
\right\}  .\label{enba}%
\end{equation}
If $\left\vert g^{N}\left(  \mathbf{n},\mathbf{\psi}_{ba}\right)
\right\rangle $ is also a stationary state then eq. $\left(  \ref{enba}%
\right)  $ becomes
\begin{equation}
\mathcal{E}\left(  \mathbf{n}_{ba},\mathbf{\psi}\right)  =\frac{1}{2}%
\sum_{1\leq p\leq2s}N_{p}I_{ba}\left(  \psi_{p}\right)  +\mathcal{E}%
_{1}\left(  \mathbf{n}_{ba},\mathbf{\psi}\right)  =\frac{1}{2}\sum_{1\leq
p\leq2s}N_{p}I_{ba}\left(  \psi_{p}\right)  +\mathcal{E}_{1}\left(
\mathbf{n},\mathbf{\psi}\right)
\end{equation}
and
\begin{align}
\mathcal{E}\left(  \mathbf{n},\mathbf{\psi}\right)  -\mathcal{E}\left(
\mathbf{n}_{ba},\mathbf{\psi}\right)   &  =N_{a}I\left(  \psi_{a}\right)
+N_{b}I\left(  \psi_{b}\right)  -N_{b}I\left(  \psi_{a}\right)  -N_{a}I\left(
\psi_{b}\right)  \nonumber\\
&  =\left(  N_{a}-N_{b}\right)  \left\{  I\left(  \psi_{a}\right)  -I\left(
\psi_{b}\right)  \right\}  >0.
\end{align}
Thus one concludes that
\begin{equation}
\mathcal{E}\left(  \mathbf{n},\mathbf{\psi}\right)  <\mathcal{E}\left(
\mathbf{n}_{ba},\mathbf{\psi}\right)
\end{equation}
which is a contradiction, hence
\begin{equation}
N_{a}>N_{b}.
\end{equation}%
\begin{align}
&  \mathcal{E}\left(  \mathbf{n}_{ba},\mathbf{\psi}\right)  \nonumber\\
&  =\frac{1}{2}%
%TCIMACRO{\tsum \limits_{1\leq p\leq2s}}%
%BeginExpansion
{\textstyle\sum\limits_{1\leq p\leq2s}}
%EndExpansion
\left\{  \frac{1}{S_{N}\left(  \mathbf{n}_{ba},\mathbf{\psi}\right)
}\left\langle \left.  g^{N}\left(  \mathbf{n}_{ba},\mathbf{\psi}\right)
\right\vert \left\{  a_{p}^{\dagger}a_{p}H-a_{p}^{\dagger}Ha_{p}\right\}
g^{N}\left(  \mathbf{n}_{ba},\mathbf{\psi}\right)  \right\rangle \right.  \\
&  +\left.  N_{p}\left\langle \left.  \psi_{p}\right\vert h_{1}\psi
_{p}\right\rangle \right\}  \\
&  =\frac{1}{2}\left\{
%TCIMACRO{\tsum \limits_{1\leq p\leq2s}}%
%BeginExpansion
{\textstyle\sum\limits_{1\leq p\leq2s}}
%EndExpansion
\frac{1}{S_{N}\left(  \mathbf{n},\mathbf{\psi}\right)  }\left\langle \left.
g^{N}\left(  \mathbf{n}_{ba},\mathbf{\psi}\right)  \right\vert \left\{
a_{p}^{\dagger}a_{p}H-a_{p}^{\dagger}Ha_{p}\right\}  g^{N}\left(
\mathbf{n}_{ba},\mathbf{\psi}\right)  \right\rangle +N_{p}\left\langle \left.
\psi_{p}\right\vert h_{1}\psi_{p}\right\rangle \right\}
\end{align}%
\begin{align}
\mathcal{E}\left(  \mathbf{n}_{ba},\mathbf{\psi}\right)   &  =\frac{1}%
{2}\left\{  \frac{2N}{S_{N}\left(  \mathbf{n},\mathbf{\psi}\right)
}\left\langle \left.  g^{N}\left(  \mathbf{n}_{ba},\mathbf{\psi}\right)
\right\vert Hg^{N}\left(  \mathbf{n}_{ba},\mathbf{\psi}\right)  \right\rangle
-%
%TCIMACRO{\tsum \limits_{1\leq p\leq2s}}%
%BeginExpansion
{\textstyle\sum\limits_{1\leq p\leq2s}}
%EndExpansion
\left\langle \left.  g^{N}\left(  \mathbf{n}_{ba},\mathbf{\psi}\right)
\right\vert a_{p}^{\dagger}Ha_{p}g^{N}\left(  \mathbf{n}_{ba},\mathbf{\psi
}\right)  \right\rangle \right\}  \\
&  \qquad\qquad\qquad\qquad\qquad\qquad\qquad\qquad\qquad\qquad\qquad
\qquad\qquad+\frac{1}{2}\mathcal{E}_{1}\left(  \mathbf{n},\mathbf{\psi
}\right)
\end{align}%
\begin{equation}
F\left(  \mathbb{D}^{2},\mathbf{\varphi}\right)  _{pp}=\frac{1}%
{\operatorname{Tr}\left\{  D^{2N}\right\}  }\operatorname{Tr}\left\{
\overset{}{\left\{  u^{\dagger}a_{p}^{\dagger}uu^{\dagger}a_{p}uu^{\dagger
}Hu-u^{\dagger}a_{p}uu^{\dagger}Huu^{\dagger}a_{p}^{\dagger}u\right\}  }%
D^{2N}\right\}
\end{equation}%
\begin{equation}
\operatorname{Tr}\left\{  \left(  \mathbb{H}_{1}\left(  \mathbf{\varphi
}\right)  \mathbb{D}^{1}\right)  _{pp}+\frac{1}{2}\left(  L_{2}^{1}\left(
\mathbb{V}_{2}\left(  \mathbf{\varphi}\right)  \mathbb{D}^{2}\right)  \right)
_{pp}\right\}
\end{equation}
The aufbau principle is also based on these eigenvalues i.e. the optimized
IPS\ is in general composed of the orbitals corresponding to most negative
ionization potentials.


\end{document}